\documentclass{article}
\usepackage[utf8]{inputenc}
\usepackage{amsmath}
\usepackage{geometry}
\geometry{margin=2cm}
\begin{document}

\section*{Análise da Questão 179 — ENEM 2024 — Matemática}

Uma empresa de engenharia executa um serviço orçado em R\$\,71\,250,00. Foi decidido que 40\% desse valor será pago a três engenheiros, proporcionalmente às horas trabalhadas. Sabemos as jornadas deles:
\begin{itemize}
  \item Engenheiro I: 4 dias, 5,5 horas por dia;
  \item Engenheiro II: 5 dias, 4 horas por dia;
  \item Engenheiro III: 6 dias, 2,5 horas por dia.
\end{itemize}

\section{Solução Técnica}

Valor total destinado aos engenheiros:
\[
0{,}4 \times 71\,250 = 28\,500 \text{ reais}
\]

Horas trabalhadas por cada um:
\begin{align*}
\text{Engenheiro I} &= 4 \times 5{,}5 = 22,\\
\text{Engenheiro II} &= 5 \times 4 = 20,\\
\text{Engenheiro III} &= 6 \times 2{,}5 = 15.
\end{align*}

Total de horas:
\[
22 + 20 + 15 = 57
\]

Valor pago por hora:
\[
\frac{28\,500}{57} = 500 \text{ reais/hora}
\]

Pagamentos individuais:
\begin{align*}
\text{I} &= 22 \times 500 = 11\,000, \\
\text{II} &= 20 \times 500 = 10\,000, \\
\text{III} &= 15 \times 500 = 7\,500.
\end{align*}

Diferenças entre pagamentos:
\begin{align*}
|\text{I} - \text{II}| &= 1\,000, \\
|\text{I} - \text{III}| &= 3\,500, \\
|\text{II} - \text{III}| &= 2\,500.
\end{align*}

Maior diferença:
\[
3\,500 \text{ reais entre os engenheiros I e III}.
\]

\section{Explicação Didática}

A solução requer calcular primeiro todas as horas trabalhadas para depois distribuir o valor proporcionalmente. Multiplicando jornadas pelas quantidades de dias, somamos o total de horas para definir o valor unitário por hora. Multiplicamos, então, o valor por hora pelas horas individuais para determinar os ganhos individuais, e finalmente comparamos para achar a maior diferença.

\section{Considerações Finais}

A resposta correta é \textbf{R\$\,3\,500}, correspondente à maior diferença entre os pagamentos de dois engenheiros. Os erros comuns envolvem confundir diferenças mínima e máxima, cálculo incorreto das horas e arredondamentos errados.

Este problema ajuda a exercitar raciocínio proporcional, interpretação e cálculo com porcentagem, importantes para a compreensão de situações reais envolvendo repartição de valores proporcionalmente ao trabalho efetuado.

\end{document}